\documentclass[12pt,a4paper]{article}

\usepackage[utf8]{inputenc}
\usepackage[T1]{fontenc}
\usepackage[spanish]{babel}
\usepackage{geometry}
\usepackage{longtable}
\usepackage{array}
\usepackage[table]{xcolor}
\usepackage{hyperref}
\usepackage{enumitem}
\usepackage{fancyhdr}
\usepackage{titlesec}

\geometry{margin=2.5cm}
\hypersetup{colorlinks=true, linkcolor=blue, urlcolor=blue}
\setlength{\parindent}{0pt}
\setlength{\parskip}{6pt}
\renewcommand{\arraystretch}{1.25}

\definecolor{azulTekMess}{HTML}{1F4E78}
\definecolor{grisClaro}{HTML}{F2F2F2}

\pagestyle{fancy}
\fancyhf{}
\lhead{Plan de Pruebas Funcionales}
\rhead{SNAAR - TekMess}
\cfoot{\thepage}

\titleformat{\section}
  {\normalfont\Large\bfseries\color{azulTekMess}}
  {\thesection}{1em}{}
\titleformat{\subsection}
  {\normalfont\large\bfseries\color{azulTekMess}}
  {\thesubsection}{1em}{}

\title{\textbf{Plan de Pruebas Funcionales}\\Sistema SNAAR - TekMess}
\author{Danna Andrade \and Ariel Llumiquinga \and David Pilaguano}
\date{9 de julio de 2026}

\begin{document}

\maketitle
\thispagestyle{empty}

\begin{center}
\begin{tabular}{|p{4.5cm}|p{8.5cm}|}
\hline
\rowcolor{grisClaro}
\textbf{Documento} & Plan de pruebas funcionales del prototipo TekMess \\
\hline
\textbf{Versión} & 1.0 \\
\hline
\textbf{Fecha de elaboración} & 9 de julio de 2026 \\
\hline
\textbf{Proyecto} & Sistema SNAAR - TekMess \\
\hline
\textbf{Tipo de pruebas} & Pruebas funcionales automatizadas con JUnit \\
\hline
\textbf{Responsables} & Danna Andrade, Ariel Llumiquinga, David Pilaguano \\
\hline
\end{tabular}
\end{center}

\newpage
\tableofcontents
\newpage

\section{Control del Documento}

\subsection{Historial de Versiones}

\begin{longtable}{|p{2cm}|p{3cm}|p{5.5cm}|p{3.5cm}|}
\hline
\rowcolor{grisClaro}
\textbf{Versión} & \textbf{Fecha} & \textbf{Descripción} & \textbf{Responsable} \\
\hline
1.0 & 09/07/2026 & Elaboración inicial del plan de pruebas funcionales, matriz de casos y trazabilidad de requisitos. & Equipo TekMess \\
\hline
\end{longtable}

\subsection{Propósito del Documento}

El propósito de este documento es establecer una guía formal para planificar, diseñar y ejecutar las pruebas funcionales del prototipo del sistema SNAAR - TekMess. El plan define los objetivos, alcance, estrategia, ambiente, criterios, datos de prueba, responsabilidades y casos que permiten verificar el cumplimiento de los requisitos funcionales identificados.

Este documento también sirve como evidencia técnica del proceso de aseguramiento de calidad aplicado al prototipo, permitiendo relacionar cada requisito con al menos un caso de prueba automatizado o verificable.

\section{Introducción}

El sistema SNAAR - TekMess es un prototipo web desarrollado en Java orientado a la gestión de empleados, generación de credenciales, autenticación por rol y administración de reportes analíticos. Debido a que el sistema maneja información de usuarios, roles y accesos, resulta necesario validar que las reglas de negocio se comporten de manera correcta antes de avanzar hacia pruebas integrales con base de datos y servidor de aplicaciones.

Las pruebas propuestas se enfocan en las funcionalidades existentes en el prototipo y se desarrollan con JUnit 5. Para aislar la lógica de negocio y evitar dependencia directa de PostgreSQL durante la ejecución unitaria, se emplean objetos DAO en memoria que simulan el comportamiento esperado de la persistencia.

La planificación considera los requisitos REQ001 a REQ009, relacionados con registro de empleados, listado, edición, eliminación, inicio de sesión, cambio de contraseña, recuperación de contraseña, generación de reportes e historial de reportes. Las pruebas fueron diseñadas desde cero con base en los requisitos entregados, sin reutilizar pruebas previas que pudieran existir en documentos externos.

\section{Objetivos}

\subsection{Objetivo General}

Validar de forma sistemática y automatizada que el prototipo TekMess cumpla los requisitos funcionales REQ001 a REQ009, verificando los flujos principales, reglas de validación, estados del sistema y respuestas esperadas en los módulos de empleados, autenticación, contraseñas y reportes.

\subsection{Objetivos Específicos}

\begin{itemize}[leftmargin=*]
    \item Verificar que el sistema permita registrar empleados válidos y genere credenciales de acceso iniciales.
    \item Comprobar que la validación de datos impida registros con cédulas inválidas, roles incorrectos, correos no válidos o datos duplicados.
    \item Confirmar que el listado de personal presente la información principal de cada empleado registrado.
    \item Validar la edición de datos de empleados existentes, manteniendo reglas de integridad y consistencia.
    \item Verificar la eliminación de empleados y la baja del usuario asociado.
    \item Evaluar el inicio de sesión por rol, el control de intentos fallidos y el bloqueo de cuentas.
    \item Comprobar el cambio de contraseña con política de seguridad, cifrado y actualización de primer acceso.
    \item Validar la recuperación de contraseña mediante usuario y correo registrados.
    \item Verificar la generación de reportes analíticos y la consulta del historial con anotaciones.
    \item Registrar la trazabilidad entre requisitos, casos de prueba, datos utilizados y resultados esperados.
\end{itemize}

\section{Datos Generales del Proyecto}

\begin{longtable}{|p{5cm}|p{9cm}|}
\hline
\rowcolor{grisClaro}
\textbf{Campo} & \textbf{Descripción} \\
\hline
Nombre del proyecto & Sistema SNAAR - TekMess \\
\hline
Tipo de producto & Prototipo web Java basado en arquitectura MVC, servicios, DAOs y vistas JSP \\
\hline
Lenguaje principal & Java 17 \\
\hline
Herramienta de construcción & Maven \\
\hline
Framework de pruebas & JUnit 5 \\
\hline
Plugin de ejecución & Maven Surefire \\
\hline
Base de datos prevista & PostgreSQL \\
\hline
Módulos evaluados & Gestión de empleados, autenticación, cambio y recuperación de contraseña, reportes analíticos \\
\hline
Artefactos generados & Pruebas automatizadas JUnit, documento LaTeX del plan de pruebas y tablero Kanban en Excel \\
\hline
Fecha del plan & 9 de julio de 2026 \\
\hline
\end{longtable}

\section{Involucrados}

\begin{longtable}{|p{4cm}|p{4cm}|p{6cm}|}
\hline
\rowcolor{grisClaro}
\textbf{Integrante} & \textbf{Rol en el proyecto} & \textbf{Responsabilidad en pruebas} \\
\hline
Danna Andrade & Integrante del equipo de desarrollo & Revisión de requisitos de gestión de empleados, validación de registro, eliminación y documentación del tablero Kanban. \\
\hline
Ariel Llumiquinga & Integrante del equipo de desarrollo & Revisión de requisitos de autenticación, inicio de sesión, bloqueo por intentos y trazabilidad de casos de prueba. \\
\hline
David Pilaguano & Integrante del equipo de desarrollo & Revisión de requisitos de reportes, historial, anotaciones y consistencia técnica del documento. \\
\hline
\end{longtable}

\section{Alcance}

\subsection{Funcionalidades Incluidas}

El plan cubre las siguientes funcionalidades del prototipo:

\begin{itemize}[leftmargin=*]
    \item Registro de empleados y generación automática de credenciales.
    \item Listado de empleados registrados.
    \item Edición de información de empleados.
    \item Eliminación de empleados y usuario asociado.
    \item Inicio de sesión por rol y bloqueo por intentos fallidos.
    \item Cambio de contraseña con política de seguridad.
    \item Recuperación de contraseña mediante usuario y correo.
    \item Generación de reportes analíticos.
    \item Consulta de historial de reportes y registro de anotaciones.
\end{itemize}

\subsection{Funcionalidades No Incluidas}

Quedan fuera del alcance de este plan:

\begin{itemize}[leftmargin=*]
    \item Pruebas de rendimiento, carga o estrés.
    \item Pruebas de seguridad avanzadas como inyección SQL, análisis de vulnerabilidades o pruebas de penetración.
    \item Pruebas visuales completas sobre JSP, CSS o experiencia de usuario.
    \item Pruebas de compatibilidad entre navegadores.
    \item Pruebas de integración real contra PostgreSQL, salvo que posteriormente se configure un entorno de base de datos controlado.
    \item Validación de exportación PDF con inspección visual del archivo generado.
\end{itemize}

\section{Estrategia de Pruebas}

La estrategia se basa en pruebas funcionales automatizadas sobre la capa de servicios y utilidades. Esta decisión permite validar las reglas de negocio de forma rápida, repetible y controlada, sin depender de elementos externos como servidor web, sesiones HTTP o base de datos real.

\begin{longtable}{|p{4cm}|p{10cm}|}
\hline
\rowcolor{grisClaro}
\textbf{Elemento} & \textbf{Estrategia Aplicada} \\
\hline
Nivel de prueba & Prueba funcional unitaria sobre servicios Java. \\
\hline
Enfoque & Validación de entradas, reglas de negocio, mensajes de respuesta y cambios de estado. \\
\hline
Aislamiento & Uso de DAOs en memoria para simular empleados, usuarios y reportes. \\
\hline
Automatización & Casos implementados en JUnit 5 dentro de \texttt{src/test/java}. \\
\hline
Datos de prueba & Datos controlados para empleados, usuarios, roles, contraseñas y reportes. \\
\hline
Ejecución & Comando esperado: \texttt{mvn test}. \\
\hline
Evidencia & Resultado de ejecución JUnit, código fuente de pruebas y trazabilidad incluida en este documento. \\
\hline
\end{longtable}

\section{Ambiente de Pruebas}

\begin{longtable}{|p{5cm}|p{9cm}|}
\hline
\rowcolor{grisClaro}
\textbf{Recurso} & \textbf{Especificación} \\
\hline
Sistema operativo & Windows o entorno compatible con Java y Maven \\
\hline
JDK & Java 17 o superior \\
\hline
Gestor de dependencias & Maven configurado en PATH o Maven Wrapper del proyecto \\
\hline
Framework de pruebas & JUnit Jupiter 5.10.2 \\
\hline
Plugin de pruebas & Maven Surefire 3.2.5 \\
\hline
Base de datos & No requerida para la ejecución unitaria; se simula mediante DAOs en memoria \\
\hline
Comando de ejecución & \texttt{mvn test} \\
\hline
\end{longtable}

\section{Criterios de Entrada y Salida}

\subsection{Criterios de Entrada}

\begin{itemize}[leftmargin=*]
    \item El código fuente del prototipo debe estar disponible.
    \item El archivo \texttt{pom.xml} debe incluir JUnit 5 y Maven Surefire.
    \item Los requisitos funcionales REQ001 a REQ009 deben estar identificados.
    \item Las clases de servicio y entidades deben compilar correctamente.
    \item El entorno de ejecución debe contar con Java y Maven.
\end{itemize}

\subsection{Criterios de Salida}

\begin{itemize}[leftmargin=*]
    \item Cada requisito funcional debe tener al menos un caso de prueba asociado.
    \item Las pruebas automatizadas deben ejecutarse sin errores.
    \item Los resultados esperados deben coincidir con los resultados obtenidos.
    \item Las incidencias detectadas deben registrarse para corrección.
    \item El plan de pruebas y el tablero Kanban deben quedar actualizados como evidencia del trabajo realizado.
\end{itemize}

\section{Datos de Prueba}

\begin{longtable}{|p{3.5cm}|p{4.5cm}|p{6cm}|}
\hline
\rowcolor{grisClaro}
\textbf{Dato} & \textbf{Valor de Referencia} & \textbf{Uso Principal} \\
\hline
Cédula válida & 1712345678 & Registro, edición, eliminación y recuperación de usuario. \\
\hline
Nombre de empleado & Danna Andrade & Generación de usuario base \texttt{dandrade}. \\
\hline
Correo institucional & danna@tekmess.com & Validación de correo y recuperación de contraseña. \\
\hline
Rol válido & SUPERVISOR, GUARDIA, JEFE\_LOGISTICA & Validación de permisos y creación de sesión por rol. \\
\hline
Contraseña temporal & Temp123! & Inicio de sesión y cambio de contraseña. \\
\hline
Contraseña nueva & Nueva123! / Recupera123! & Cambio y recuperación de contraseña cumpliendo política. \\
\hline
Rango de fechas válido & Fecha inicio menor que fecha fin & Generación y consulta de reportes. \\
\hline
\end{longtable}

\section{Matriz de Trazabilidad}

\begin{longtable}{|p{2cm}|p{5cm}|p{4cm}|p{3cm}|}
\hline
\rowcolor{grisClaro}
\textbf{Requisito} & \textbf{Funcionalidad Evaluada} & \textbf{Clase Principal} & \textbf{Casos} \\
\hline
REQ001 & Registro de empleados y generación de credenciales & \texttt{EmpleadoServicio}, \texttt{GeneradorCredenciales}, \texttt{ValidadorDatos} & TC-001, TC-002 \\
\hline
REQ002 & Listado de personal registrado & \texttt{EmpleadoServicio} & TC-003 \\
\hline
REQ003 & Edición de datos del empleado & \texttt{EmpleadoServicio} & TC-004 \\
\hline
REQ004 & Eliminación de empleados & \texttt{EmpleadoServicio} & TC-005 \\
\hline
REQ005 & Inicio de sesión por rol & \texttt{AuthServicio}, \texttt{Sesion} & TC-006, TC-007 \\
\hline
REQ006 & Cambio de contraseña & \texttt{AuthServicio}, \texttt{CifradorContrasena} & TC-008 \\
\hline
REQ007 & Recuperación de contraseña & \texttt{AuthServicio} & TC-009 \\
\hline
REQ008 & Generación de reporte analítico & \texttt{ReporteServicio} & TC-010 \\
\hline
REQ009 & Consulta de historial de reportes & \texttt{ReporteServicio} & TC-011, TC-012 \\
\hline
\end{longtable}

\section{Desglose Detallado de Pruebas Funcionales}

En esta sección se describe cada prueba de forma individual. Para cada caso se especifica el objetivo, requisito funcional relacionado, precondiciones, datos utilizados, procedimiento, elementos evaluados, resultado esperado y criterio de aprobación. Esta estructura permite comprender exactamente qué se valida en cada funcionalidad sin reducir el análisis a una matriz resumida.

\subsection{TC-001: Registro correcto de empleado y generación de credenciales}

\textbf{Requisito asociado:} REQ001 - Registro de empleados y generación de credenciales.

\textbf{Objetivo de la prueba:} comprobar que el sistema permita registrar un empleado con datos válidos y que, como consecuencia del registro, genere automáticamente una cuenta de usuario con credenciales temporales.

\textbf{Precondiciones:}
\begin{itemize}[leftmargin=*]
    \item No debe existir un empleado registrado con la cédula utilizada en la prueba.
    \item El DAO de empleados debe estar disponible en memoria.
    \item El DAO de usuarios debe estar disponible en memoria.
\end{itemize}

\textbf{Datos de prueba:}
\begin{itemize}[leftmargin=*]
    \item Cédula: \texttt{1712345678}.
    \item Nombres: \texttt{Danna Andrade}.
    \item Correo: \texttt{DANNA.ANDRADE@TEKMESS.COM}.
    \item Rol: \texttt{JEFE\_LOGISTICA}.
\end{itemize}

\textbf{Procedimiento:}
\begin{enumerate}[leftmargin=*]
    \item Crear una instancia de \texttt{EmpleadoServicio} con DAOs en memoria.
    \item Construir un objeto \texttt{Empleado} con datos válidos.
    \item Ejecutar el método \texttt{crearEmpleado}.
    \item Consultar el empleado almacenado por cédula.
    \item Consultar el usuario generado para la cédula registrada.
\end{enumerate}

\textbf{Aspectos evaluados:}
\begin{itemize}[leftmargin=*]
    \item Validación de cédula con exactamente diez dígitos numéricos.
    \item Validación de nombres completos.
    \item Validación de rol permitido por el sistema.
    \item Normalización del correo institucional a minúsculas.
    \item Generación de nombre de usuario a partir de nombres y apellidos.
    \item Generación de contraseña temporal que cumpla política de seguridad.
    \item Cifrado de la contraseña mediante BCrypt.
    \item Marcación del usuario con primer acceso obligatorio.
\end{itemize}

\textbf{Resultado esperado:} el servicio debe retornar un mensaje de registro exitoso, el empleado debe existir en el DAO, el usuario debe estar asociado a la cédula del empleado, la contraseña temporal debe cumplir la política de seguridad y el campo de primer acceso debe quedar activo.

\textbf{Criterio de aprobación:} la prueba se aprueba si todas las validaciones anteriores se cumplen y la contraseña temporal coincide con el hash almacenado al verificarla con el cifrador.

\subsection{TC-002: Rechazo de registro con cédula duplicada}

\textbf{Requisito asociado:} REQ001 - Registro de empleados y generación de credenciales.

\textbf{Objetivo de la prueba:} verificar que el sistema impida registrar un nuevo empleado cuando la cédula ya existe en el repositorio de empleados.

\textbf{Precondiciones:}
\begin{itemize}[leftmargin=*]
    \item Debe existir previamente un empleado registrado con la misma cédula.
    \item El servicio debe consultar la existencia de la cédula antes de crear el nuevo registro.
\end{itemize}

\textbf{Datos de prueba:}
\begin{itemize}[leftmargin=*]
    \item Cédula duplicada: \texttt{1712345678}.
    \item Primer empleado: \texttt{Danna Andrade}.
    \item Segundo empleado: \texttt{Ariel Llumiquinga}.
\end{itemize}

\textbf{Procedimiento:}
\begin{enumerate}[leftmargin=*]
    \item Registrar un empleado inicial en el DAO en memoria.
    \item Intentar registrar un segundo empleado usando la misma cédula.
    \item Capturar el mensaje retornado por \texttt{crearEmpleado}.
\end{enumerate}

\textbf{Aspectos evaluados:}
\begin{itemize}[leftmargin=*]
    \item Verificación de unicidad de la cédula.
    \item Prevención de duplicidad en el registro de personal.
    \item Control de flujo para evitar creación de usuario cuando el empleado no debe registrarse.
    \item Mensaje funcional comprensible para el usuario.
\end{itemize}

\textbf{Resultado esperado:} el sistema debe rechazar el registro y retornar un mensaje indicando que la cédula ya se encuentra registrada.

\textbf{Criterio de aprobación:} la prueba se aprueba si no se crea un nuevo empleado, no se generan credenciales adicionales y el mensaje de error corresponde al escenario de duplicidad.

\subsection{TC-003: Listado de personal registrado}

\textbf{Requisito asociado:} REQ002 - Listado de personal registrado.

\textbf{Objetivo de la prueba:} confirmar que el servicio permita consultar el conjunto de empleados registrados y que la información principal de cada empleado esté disponible para la vista de listado.

\textbf{Precondiciones:}
\begin{itemize}[leftmargin=*]
    \item Deben existir empleados cargados en el DAO en memoria.
    \item Cada empleado debe tener cédula, nombres, correo y rol.
\end{itemize}

\textbf{Datos de prueba:}
\begin{itemize}[leftmargin=*]
    \item Empleado 1: \texttt{Danna Andrade}, rol \texttt{JEFE\_LOGISTICA}.
    \item Empleado 2: \texttt{Ariel Llumiquinga}, rol \texttt{SUPERVISOR}.
\end{itemize}

\textbf{Procedimiento:}
\begin{enumerate}[leftmargin=*]
    \item Registrar dos empleados en el DAO en memoria.
    \item Ejecutar el método \texttt{consultarEmpleados}.
    \item Verificar la cantidad de elementos retornados.
    \item Validar que uno de los empleados contenga cédula, rol y correo esperados.
\end{enumerate}

\textbf{Aspectos evaluados:}
\begin{itemize}[leftmargin=*]
    \item Recuperación del listado completo de empleados.
    \item Conservación de datos principales del empleado.
    \item Disponibilidad de información necesaria para edición o eliminación posterior.
\end{itemize}

\textbf{Resultado esperado:} el método debe retornar una lista con los empleados registrados y sus datos principales.

\textbf{Criterio de aprobación:} la prueba se aprueba si la lista contiene la cantidad esperada de empleados y los datos consultados coinciden con los registros cargados.

\subsection{TC-004: Edición de datos de un empleado existente}

\textbf{Requisito asociado:} REQ003 - Edición de datos del empleado.

\textbf{Objetivo de la prueba:} validar que el sistema permita modificar datos editables de un empleado existente, conservando la cédula como identificador principal.

\textbf{Precondiciones:}
\begin{itemize}[leftmargin=*]
    \item Debe existir un empleado registrado con la cédula indicada.
    \item Los nuevos datos deben cumplir las reglas de validación.
\end{itemize}

\textbf{Datos de prueba:}
\begin{itemize}[leftmargin=*]
    \item Cédula: \texttt{1712345678}.
    \item Correo inicial: \texttt{danna@tekmess.com}.
    \item Correo actualizado: \texttt{danna.andrade@tekmess.com}.
    \item Rol actualizado: \texttt{SUPERVISOR}.
\end{itemize}

\textbf{Procedimiento:}
\begin{enumerate}[leftmargin=*]
    \item Crear un empleado inicial en el DAO.
    \item Construir un objeto \texttt{Empleado} con la misma cédula y datos modificados.
    \item Ejecutar el método \texttt{editarEmpleado}.
    \item Consultar el empleado por cédula.
    \item Comparar los datos actualizados.
\end{enumerate}

\textbf{Aspectos evaluados:}
\begin{itemize}[leftmargin=*]
    \item Búsqueda de empleado existente por cédula.
    \item Validación de correo institucional.
    \item Validación de rol permitido.
    \item Persistencia de cambios en el repositorio.
    \item Mensaje de confirmación de actualización.
\end{itemize}

\textbf{Resultado esperado:} el empleado debe quedar actualizado con el nuevo correo y rol, y el servicio debe informar que los datos fueron actualizados exitosamente.

\textbf{Criterio de aprobación:} la prueba se aprueba si el registro consultado después de la edición refleja los cambios enviados y el contador de ediciones del DAO en memoria aumenta.

\subsection{TC-005: Eliminación de empleado y usuario asociado}

\textbf{Requisito asociado:} REQ004 - Eliminación de empleados.

\textbf{Objetivo de la prueba:} verificar que al eliminar un empleado también se elimine o invalide la cuenta de usuario asociada a su cédula.

\textbf{Precondiciones:}
\begin{itemize}[leftmargin=*]
    \item Debe existir un empleado registrado.
    \item Debe existir un usuario asociado a la misma cédula.
\end{itemize}

\textbf{Datos de prueba:}
\begin{itemize}[leftmargin=*]
    \item Cédula: \texttt{1712345678}.
    \item Usuario asociado: \texttt{dandrade}.
    \item Contraseña temporal inicial: \texttt{Temp123!}.
\end{itemize}

\textbf{Procedimiento:}
\begin{enumerate}[leftmargin=*]
    \item Registrar un empleado en el DAO de empleados.
    \item Registrar un usuario asociado en el DAO de usuarios.
    \item Ejecutar el método \texttt{eliminarEmpleado}.
    \item Consultar nuevamente el empleado por cédula.
    \item Consultar nuevamente el usuario por cédula.
\end{enumerate}

\textbf{Aspectos evaluados:}
\begin{itemize}[leftmargin=*]
    \item Localización del empleado antes de eliminar.
    \item Eliminación del usuario asociado.
    \item Eliminación del registro de empleado.
    \item Mensaje de eliminación exitosa.
    \item Prevención de accesos posteriores con usuario de empleado eliminado.
\end{itemize}

\textbf{Resultado esperado:} el empleado y el usuario asociado no deben existir después de ejecutar la eliminación.

\textbf{Criterio de aprobación:} la prueba se aprueba si ambas consultas posteriores retornan \texttt{null} y el servicio informa eliminación exitosa.

\subsection{TC-006: Inicio de sesión válido por rol}

\textbf{Requisito asociado:} REQ005 - Inicio de sesión por rol.

\textbf{Objetivo de la prueba:} comprobar que un usuario activo con credenciales correctas pueda iniciar sesión y que el sistema construya una sesión con el rol correspondiente.

\textbf{Precondiciones:}
\begin{itemize}[leftmargin=*]
    \item Debe existir un empleado con rol definido.
    \item Debe existir un usuario activo vinculado a la cédula del empleado.
    \item La contraseña ingresada debe coincidir con el hash almacenado.
\end{itemize}

\textbf{Datos de prueba:}
\begin{itemize}[leftmargin=*]
    \item Usuario: \texttt{dandrade}.
    \item Contraseña: \texttt{Temp123!}.
    \item Rol esperado: \texttt{SUPERVISOR}.
\end{itemize}

\textbf{Procedimiento:}
\begin{enumerate}[leftmargin=*]
    \item Crear un empleado con rol \texttt{SUPERVISOR}.
    \item Crear un usuario activo asociado al empleado.
    \item Ejecutar \texttt{iniciarSesion} con usuario y contraseña correctos.
    \item Verificar que el objeto retornado sea una instancia de \texttt{Sesion}.
    \item Verificar el rol asignado a la sesión.
\end{enumerate}

\textbf{Aspectos evaluados:}
\begin{itemize}[leftmargin=*]
    \item Validación de formato del nombre de usuario.
    \item Verificación de contraseña cifrada.
    \item Revisión de estado de cuenta activo.
    \item Reinicio de intentos fallidos tras autenticación correcta.
    \item Creación de sesión con rol funcional.
    \item Redirección a cambio de contraseña cuando el usuario está en primer acceso.
\end{itemize}

\textbf{Resultado esperado:} el sistema debe retornar una sesión válida con rol \texttt{SUPERVISOR} y el mensaje \texttt{PRIMER\_ACCESO} cuando corresponda.

\textbf{Criterio de aprobación:} la prueba se aprueba si la sesión no es nula, el rol coincide con el empleado asociado y la respuesta funcional corresponde al estado de primer acceso.

\subsection{TC-007: Bloqueo de cuenta por intentos fallidos}

\textbf{Requisito asociado:} REQ005 - Inicio de sesión por rol.

\textbf{Objetivo de la prueba:} verificar que el sistema controle los intentos fallidos de autenticación y bloquee la cuenta al alcanzar el límite permitido.

\textbf{Precondiciones:}
\begin{itemize}[leftmargin=*]
    \item Debe existir un usuario activo.
    \item La cuenta no debe estar bloqueada al inicio de la prueba.
\end{itemize}

\textbf{Datos de prueba:}
\begin{itemize}[leftmargin=*]
    \item Usuario: \texttt{dandrade}.
    \item Contraseña incorrecta: \texttt{Mala123!}.
    \item Límite de intentos: 3.
\end{itemize}

\textbf{Procedimiento:}
\begin{enumerate}[leftmargin=*]
    \item Ejecutar \texttt{iniciarSesion} con contraseña incorrecta.
    \item Repetir la operación hasta completar tres intentos fallidos.
    \item Consultar el estado de la cuenta.
    \item Revisar el mensaje funcional retornado por el tercer intento.
\end{enumerate}

\textbf{Aspectos evaluados:}
\begin{itemize}[leftmargin=*]
    \item Incremento del contador de intentos fallidos.
    \item Persistencia del número de intentos.
    \item Cambio de estado de cuenta a \texttt{BLOQUEADO}.
    \item Respuesta funcional ante cuenta bloqueada.
    \item Prevención de acceso con credenciales incorrectas.
\end{itemize}

\textbf{Resultado esperado:} luego del tercer intento fallido, el usuario debe quedar con estado \texttt{BLOQUEADO} y el sistema debe retornar un mensaje de bloqueo.

\textbf{Criterio de aprobación:} la prueba se aprueba si no se crea sesión, el estado de la cuenta cambia a bloqueado y el mensaje informa que se debe contactar al administrador.

\subsection{TC-008: Cambio de contraseña válido}

\textbf{Requisito asociado:} REQ006 - Cambio de contraseña.

\textbf{Objetivo de la prueba:} comprobar que un usuario pueda cambiar su contraseña cuando ingresa correctamente la clave actual y la nueva contraseña cumple la política de seguridad.

\textbf{Precondiciones:}
\begin{itemize}[leftmargin=*]
    \item Debe existir un usuario registrado.
    \item La contraseña actual debe estar cifrada.
    \item El usuario debe estar marcado inicialmente con primer acceso obligatorio.
\end{itemize}

\textbf{Datos de prueba:}
\begin{itemize}[leftmargin=*]
    \item Usuario: \texttt{dandrade}.
    \item Contraseña actual: \texttt{Temp123!}.
    \item Nueva contraseña: \texttt{Nueva123!}.
    \item Confirmación: \texttt{Nueva123!}.
\end{itemize}

\textbf{Procedimiento:}
\begin{enumerate}[leftmargin=*]
    \item Crear un usuario con contraseña temporal cifrada.
    \item Ejecutar \texttt{cambiarContrasena}.
    \item Consultar el usuario actualizado.
    \item Verificar que el nuevo hash corresponda a la nueva contraseña.
    \item Verificar que el indicador de primer acceso quede desactivado.
\end{enumerate}

\textbf{Aspectos evaluados:}
\begin{itemize}[leftmargin=*]
    \item Obligatoriedad de contraseña actual.
    \item Coincidencia entre nueva contraseña y confirmación.
    \item Cumplimiento de política: longitud mínima, mayúscula, minúscula, dos números y carácter especial.
    \item Rechazo de contraseña igual a la actual.
    \item Actualización segura mediante hash.
    \item Desactivación del primer acceso obligatorio.
\end{itemize}

\textbf{Resultado esperado:} el sistema debe actualizar la contraseña, almacenar un nuevo hash y retornar un mensaje de actualización exitosa.

\textbf{Criterio de aprobación:} la prueba se aprueba si el nuevo hash valida la contraseña nueva y el usuario deja de estar marcado como primer acceso.

\subsection{TC-009: Recuperación de contraseña con correo y usuario registrados}

\textbf{Requisito asociado:} REQ007 - Recuperación de contraseña.

\textbf{Objetivo de la prueba:} validar que un usuario pueda restablecer su contraseña cuando proporciona un nombre de usuario y correo electrónico registrados en el sistema.

\textbf{Precondiciones:}
\begin{itemize}[leftmargin=*]
    \item Debe existir un empleado con correo registrado.
    \item Debe existir un usuario asociado a la cédula del empleado.
    \item La cuenta no debe estar bloqueada.
\end{itemize}

\textbf{Datos de prueba:}
\begin{itemize}[leftmargin=*]
    \item Correo: \texttt{danna@tekmess.com}.
    \item Usuario: \texttt{dandrade}.
    \item Nueva contraseña: \texttt{Recupera123!}.
    \item Confirmación: \texttt{Recupera123!}.
\end{itemize}

\textbf{Procedimiento:}
\begin{enumerate}[leftmargin=*]
    \item Registrar empleado y usuario asociado.
    \item Ejecutar \texttt{recuperarContrasena}.
    \item Consultar el usuario por nombre de usuario.
    \item Verificar que el hash corresponda a la nueva contraseña.
\end{enumerate}

\textbf{Aspectos evaluados:}
\begin{itemize}[leftmargin=*]
    \item Validación de formato de usuario.
    \item Validación de formato de correo.
    \item Verificación de relación entre usuario, cédula y correo del empleado.
    \item Bloqueo del restablecimiento cuando la cuenta está bloqueada.
    \item Aplicación de la política de seguridad a la nueva contraseña.
    \item Almacenamiento cifrado de la contraseña restablecida.
\end{itemize}

\textbf{Resultado esperado:} el sistema debe restablecer la contraseña y confirmar la operación mediante un mensaje exitoso.

\textbf{Criterio de aprobación:} la prueba se aprueba si el hash almacenado valida la nueva contraseña y el mensaje indica restablecimiento exitoso.

\subsection{TC-010: Generación de reporte analítico}

\textbf{Requisito asociado:} REQ008 - Generación de reporte analítico.

\textbf{Objetivo de la prueba:} verificar que el servicio consolide correctamente los datos de empleados gestionados y accesos fallidos en un reporte analítico.

\textbf{Precondiciones:}
\begin{itemize}[leftmargin=*]
    \item Debe existir un rango de fechas válido.
    \item Los DAOs deben proporcionar totales de empleados creados, editados, eliminados y accesos fallidos.
    \item Debe existir un DAO de reportes capaz de guardar el reporte generado.
\end{itemize}

\textbf{Datos de prueba:}
\begin{itemize}[leftmargin=*]
    \item Empleados creados: 2.
    \item Empleados editados: 1.
    \item Empleados eliminados: 1.
    \item Accesos fallidos: 3.
    \item Generado por: \texttt{Danna Andrade}.
\end{itemize}

\textbf{Procedimiento:}
\begin{enumerate}[leftmargin=*]
    \item Configurar los totales en los DAOs en memoria.
    \item Ejecutar \texttt{generarReporte} con fecha inicial menor que fecha final.
    \item Obtener el objeto \texttt{Reporte} retornado.
    \item Validar cada total consolidado.
    \item Confirmar el mensaje de generación exitosa.
\end{enumerate}

\textbf{Aspectos evaluados:}
\begin{itemize}[leftmargin=*]
    \item Validación de rango de fechas.
    \item Consulta de empleados creados en el período.
    \item Consulta de empleados editados en el período.
    \item Consulta de empleados eliminados en el período.
    \item Consulta de accesos fallidos.
    \item Construcción del objeto \texttt{Reporte}.
    \item Persistencia del reporte generado.
\end{itemize}

\textbf{Resultado esperado:} el servicio debe retornar un reporte con los totales esperados y el mensaje \texttt{Reporte generado exitosamente.}

\textbf{Criterio de aprobación:} la prueba se aprueba si todos los totales del reporte coinciden con los datos configurados y el reporte queda guardado.

\subsection{TC-011: Consulta de historial de reportes por rango de fechas}

\textbf{Requisito asociado:} REQ009 - Consulta de historial de reportes.

\textbf{Objetivo de la prueba:} comprobar que el sistema permita consultar reportes históricos aplicando un filtro de fechas válido.

\textbf{Precondiciones:}
\begin{itemize}[leftmargin=*]
    \item Debe existir al menos un reporte guardado.
    \item El rango de consulta debe incluir el período del reporte.
\end{itemize}

\textbf{Datos de prueba:}
\begin{itemize}[leftmargin=*]
    \item Fecha inicial de reporte: \texttt{1000 ms}.
    \item Fecha final de reporte: \texttt{2000 ms}.
    \item Rango de búsqueda: \texttt{500 ms} a \texttt{3000 ms}.
    \item Generado por: \texttt{David Pilaguano}.
\end{itemize}

\textbf{Procedimiento:}
\begin{enumerate}[leftmargin=*]
    \item Crear y guardar un reporte en el DAO en memoria.
    \item Ejecutar \texttt{consultarHistorial} con un rango válido.
    \item Verificar que el resultado no sea nulo.
    \item Validar que la lista contenga el reporte esperado.
\end{enumerate}

\textbf{Aspectos evaluados:}
\begin{itemize}[leftmargin=*]
    \item Validación de rango de fechas para consulta.
    \item Recuperación de reportes históricos.
    \item Aplicación del filtro por período.
    \item Mensaje funcional \texttt{OK} cuando existen resultados.
\end{itemize}

\textbf{Resultado esperado:} el sistema debe retornar una lista con el reporte ubicado dentro del rango de fechas y el estado de respuesta debe ser \texttt{OK}.

\textbf{Criterio de aprobación:} la prueba se aprueba si el historial contiene exactamente el reporte esperado y la respuesta funcional indica éxito.

\subsection{TC-012: Registro de anotación en un reporte existente}

\textbf{Requisito asociado:} REQ009 - Consulta de historial de reportes.

\textbf{Objetivo de la prueba:} validar que se puedan agregar anotaciones a un reporte previamente generado sin modificar los datos base del reporte.

\textbf{Precondiciones:}
\begin{itemize}[leftmargin=*]
    \item Debe existir un reporte guardado.
    \item El identificador del reporte debe ser válido.
\end{itemize}

\textbf{Datos de prueba:}
\begin{itemize}[leftmargin=*]
    \item Autor: \texttt{Ariel Llumiquinga}.
    \item Contenido de la anotación: \texttt{Revision aprobada}.
\end{itemize}

\textbf{Procedimiento:}
\begin{enumerate}[leftmargin=*]
    \item Guardar un reporte en el DAO en memoria.
    \item Ejecutar \texttt{agregarAnotacion} indicando el identificador del reporte.
    \item Consultar las anotaciones asociadas al reporte.
    \item Verificar que la anotación haya sido almacenada.
\end{enumerate}

\textbf{Aspectos evaluados:}
\begin{itemize}[leftmargin=*]
    \item Búsqueda del reporte por identificador.
    \item Creación de objeto \texttt{Anotacion}.
    \item Asociación de la anotación al reporte correcto.
    \item Conservación de los datos base del reporte.
    \item Respuesta booleana de éxito.
\end{itemize}

\textbf{Resultado esperado:} la anotación debe quedar registrada y asociada al reporte existente.

\textbf{Criterio de aprobación:} la prueba se aprueba si el método retorna \texttt{true} y la consulta de anotaciones devuelve el registro agregado.

\section{Gestión de Defectos}

Cuando una prueba falle, se deberá registrar una incidencia con la siguiente información mínima:

\begin{itemize}[leftmargin=*]
    \item Identificador del defecto.
    \item Caso de prueba asociado.
    \item Requisito funcional afectado.
    \item Descripción clara del comportamiento observado.
    \item Resultado esperado.
    \item Evidencia de ejecución.
    \item Severidad: crítica, alta, media o baja.
    \item Estado: abierto, en corrección, corregido, verificado o cerrado.
\end{itemize}

\section{Riesgos y Mitigaciones}

\begin{longtable}{|p{4.5cm}|p{5cm}|p{4.5cm}|}
\hline
\rowcolor{grisClaro}
\textbf{Riesgo} & \textbf{Impacto} & \textbf{Mitigación} \\
\hline
Maven no disponible en el entorno & No se pueden ejecutar las pruebas con \texttt{mvn test}. & Instalar Maven, configurar PATH o agregar Maven Wrapper al proyecto. \\
\hline
Dependencia de base de datos real en algunas clases & Puede dificultar pruebas unitarias puras. & Inyectar interfaces DAO y usar dobles de prueba en memoria. \\
\hline
Reglas funcionales incompletas en el prototipo & Algunos escenarios podrían requerir validación manual o integración futura. & Documentar limitaciones y ampliar pruebas conforme evolucione el código. \\
\hline
Codificación incorrecta de caracteres & Puede afectar documentación y mensajes en español. & Guardar archivos en UTF-8 y revisar salida del compilador LaTeX. \\
\hline
\end{longtable}

\section{Procedimiento de Ejecución}

\begin{enumerate}[leftmargin=*]
    \item Verificar que Java esté instalado mediante \texttt{java -version}.
    \item Verificar que Maven esté instalado mediante \texttt{mvn -version}.
    \item Ubicarse en la raíz del prototipo, donde se encuentra el archivo \texttt{pom.xml}.
    \item Ejecutar el comando \texttt{mvn test}.
    \item Revisar el resumen de ejecución generado por Maven Surefire.
    \item En caso de fallos, identificar el caso afectado y registrar la incidencia correspondiente.
\end{enumerate}

\section{Entregables}

\begin{longtable}{|p{5cm}|p{9cm}|}
\hline
\rowcolor{grisClaro}
\textbf{Entregable} & \textbf{Descripción} \\
\hline
Archivo de pruebas JUnit & Clase \texttt{RequisitosFuncionalesTest.java} con casos automatizados para REQ001 a REQ009. \\
\hline
Plan de pruebas en LaTeX & Documento \texttt{Plan\_Pruebas\_TekMess.tex} con planificación, estrategia y matriz de casos. \\
\hline
Tablero Kanban en Excel & Archivo \texttt{Tablero\_Kanban\_Pruebas\_TekMess.xlsx} con tareas ejecutadas y responsables. \\
\hline
\end{longtable}

\section{Observaciones Finales}

El plan de pruebas fortalece la verificación funcional del prototipo al cubrir cada requisito mediante casos trazables y automatizados. La estrategia aplicada permite validar reglas de negocio sin depender de una base de datos real, lo cual mejora la rapidez de ejecución y facilita la detección temprana de errores.

Durante la preparación del plan se identificó que el entorno local utilizado no contaba con el comando \texttt{mvn} disponible en el PATH ni con compilador LaTeX instalado. Por este motivo, la ejecución formal de \texttt{mvn test} y la generación del PDF deberán realizarse en un entorno que cuente con dichas herramientas configuradas. Aun así, el documento fuente LaTeX, las pruebas JUnit y el tablero Kanban quedan listos para su revisión y ejecución.

\end{document}
